\chapter{Zmienne, stale i typy}

\section{Po co istnieja zmienne}
Strona internetowa rzadko sklada sie wyłącznie ze stalych, niezmiennych napisow.
Zwykle chcemy przechowywac dane takie jak:

\begin{itemize}
  \item tytul strony,
  \item licznik wpisow,
  \item informacja, czy element ma byc pokazany,
  \item lista kart produktowych,
  \item obiekt z danymi autora albo wpisu.
\end{itemize}

Do tego wlasnie sluza zmienne. Zmienna to \important{nazwa wskazujaca na wartosc}.
W MoonChunk mozna pracowac zarowno ze zmiennymi zmienialnymi, jak i z wartosciami,
ktore po nadaniu powinny pozostac stale.

\section{\mc{let} i \mc{const}}
MoonChunk, podobnie jak wspolczesny JavaScript, rozroznia dwa podstawowe sposoby
deklarowania lokalnych nazw.

\subsection*{\mc{let}}
\mc{let} sluzy do tworzenia zmiennej, ktorej wartosc mozna pozniej zmienic.

\begin{lstlisting}[style=moonchunk,caption={Zmienna deklarowana przez let}]
let counter: int = 0;
counter = counter + 1;
\end{lstlisting}

\subsection*{\mc{const}}
\mc{const} sluzy do tworzenia stalej, czyli nazwy, ktorej nie wolno pozniej przypisac
nowej wartosci.

\begin{lstlisting}[style=moonchunk,caption={Stala deklarowana przez const}]
const title: string = "MoonChunk";
\end{lstlisting}

Jesli potem sprobowalibysmy zrobic:

\begin{lstlisting}[style=moonchunk,caption={Niedozwolona zmiana wartosci stalej}]
title = "Nowy tytul";
\end{lstlisting}

MoonChunk zglosi blad.

\section{Czy \mc{let} musi miec wartosc od razu}
Nie zawsze. Jedna z praktycznych cech MoonChunk polega na tym, ze \mc{let} moze
zostac zadeklarowany bez wartosci poczatkowej.

\begin{lstlisting}[style=moonchunk,caption={Deklaracja let bez inicjalizacji}]
let later: int;
later = 7;
\end{lstlisting}

To bywa wygodne, gdy:

\begin{itemize}
  \item jeszcze nie znasz wartosci w chwili deklaracji,
  \item wartosc zalezy od warunku,
  \item chcesz najpierw zadeklarowac nazwe, a pozniej obliczyc wynik.
\end{itemize}

Trzeba jednak uwazac: \important{taka zmienna musi zostac zainicjalizowana przed uzyciem}.
Jesli sprobowalibysmy odczytac ja za wczesnie, MoonChunk zglosi blad runtime.

\section{Czy \mc{const} moze byc puste na starcie}
Nie. \mc{const} musi od razu dostac konkretna wartosc.

\begin{lstlisting}[style=moonchunk,caption={Poprawna deklaracja const}]
const maxItems: int = 10;
\end{lstlisting}

Taka zasada ma sens: skoro nazwa ma byc stala, od razu trzeba wiedziec, jaka wartosc
ma reprezentowac.

\section{Typy w MoonChunk}
Typ okresla, jakiego rodzaju jest dana wartosc. To bardzo wazne, bo interpreter
musi wiedziec, czy pracuje z liczba, tekstem, wartoscia logiczna, tablica czy obiektem.

MoonChunk obsluguje nastepujace typy podstawowe i specjalne:

\begin{itemize}
  \item \mc{int},
  \item \mc{float},
  \item \mc{double},
  \item \mc{number},
  \item \mc{bool},
  \item \mc{string},
  \item \mc{array},
  \item \mc{object},
  \item \mc{dict},
  \item \mc{null},
  \item \mc{undefined},
  \item \mc{unknown},
  \item \mc{any},
  \item \mc{void} dla funkcji nic nie zwracajacych.
\end{itemize}

\section{Typy liczbowe}
MoonChunk rozroznia kilka wariantow liczb, bo w praktyce nie kazda liczba ma zachowywac
sie identycznie.

\subsection*{\mc{int}}
To liczba calkowita. Gdy pracujesz na samych \mc{int}, szczegolnie przy dzieleniu,
warto pamietac, ze wynik powinien pozostac liczba calkowita.

\begin{lstlisting}[style=moonchunk,caption={Dzielenie liczb typu int}]
let a: int = 1 / 2;
\end{lstlisting}

Taki zapis daje wynik \mc{0}, a nie \mc{0.5}.

\subsection*{\mc{float} i \mc{double}}
To typy liczb zmiennoprzecinkowych, przydatne wtedy, gdy chcesz zachowac czesc ulamkowa.

\begin{lstlisting}[style=moonchunk,caption={Dzielenie z czescia ulamkowa}]
let b: double = 1 / 2.0;
let c: float = 1f / 2;
\end{lstlisting}

Tutaj wynik zachowuje informacje o ulamku.

\subsection*{\mc{number}}
Mozesz traktowac \mc{number} jako bardziej ogolna kategorie liczby. Przydaje sie tam,
gdzie nie chcesz przywiazywac sie do jednego konkretnego wariantu liczbowego,
ale nadal pracujesz na liczbach.

\section{Typ \mc{bool}}
\mc{bool} przechowuje wartosc logiczna: \mc{true} albo \mc{false}. To typ idealny do:

\begin{itemize}
  \item warunkow \mc{if},
  \item sterowania petlami,
  \item flag typu \emph{czy pokazac element},
  \item porownan i sprawdzen.
\end{itemize}

\begin{lstlisting}[style=moonchunk,caption={Przyklad zmiennej logicznej}]
let isPublished: bool = true;
\end{lstlisting}

\section{Typ \mc{string}}
\mc{string} oznacza tekst. W praktyce to jeden z najczesciej uzywanych typow, bo strony
internetowe sa w duzej mierze zbudowane wlasnie z tekstu.

\begin{lstlisting}[style=moonchunk,caption={Przyklad tekstu}]
let heading: string = "Dokumentacja MoonChunk";
\end{lstlisting}

\section{Typy zlozone: \mc{array}, \mc{object}, \mc{dict}}
\subsection*{\mc{array}}
Tablica sluzy do przechowywania kolekcji elementow w okreslonej kolejnosci.

\begin{lstlisting}[style=moonchunk,caption={Tablica}]
const tags: array = ["dsl", "html", "static-site"];
\end{lstlisting}

\subsection*{\mc{object}}
Obiekt sluzy do przechowywania nazwanych pol.

\begin{lstlisting}[style=moonchunk,caption={Obiekt}]
const author: object = { name: "Alicja", role: "writer" };
\end{lstlisting}

\subsection*{\mc{dict}}
\mc{dict} to praktycznie slownikowy wariant obiektu. W MoonChunk traktowany jest jako
obiekt o strukturze klucz--wartosc. W dokumentacji uzytkowej warto myslec o nim jako o
obiekcie przeznaczonym szczegolnie do przechowywania danych.

\begin{lstlisting}[style=moonchunk,caption={Slownik typu dict}]
let point: dict = { x: 1, y: 2 };
\end{lstlisting}

\section{Typy specjalne: \mc{null} i \mc{undefined}}
Te dwa typy sa podobne tylko pozornie.

\subsection*{\mc{null}}
\mc{null} oznacza zwykle swiadomie ustawiony brak konkretnej wartosci.

\subsection*{\mc{undefined}}
\mc{undefined} oznacza raczej brak przypisanej wartosci albo stan pusty wynikajacy
z kontekstu.

W praktyce jesli dopiero zaczynasz, nie musisz naduzywac tych typow. Warto jednak
wiedziec, ze istnieja, bo pojawiaja sie w bardziej zaawansowanych scenariuszach.

\section{Typy \mc{unknown} i \mc{any}}
To typy bardziej zaawansowane. Najlatwiej myslec o nich jak o tymczasowych pudelkach,
ktore mowia interpreterowi: \emph{na chwile odloz scisle sprawdzanie typu}.

\subsection*{\mc{unknown}}
Jest bezpieczniejszy koncepcyjnie. Mowi: \emph{mam jakas wartosc, ale jeszcze nie chce
zakladac, czym ona jest}.

\subsection*{\mc{any}}
Jest bardziej permissive. Mowi: \emph{traktuj te wartosc bardzo elastycznie}.

W praktyce te typy sa najczesciej widoczne przy bardziej zaawansowanych rzutowaniach,
zwlaszcza kiedy chcesz przejsc przez zapis typu \mc{as unknown as ...}.

\section{Typ \mc{void}}
\mc{void} dotyczy przede wszystkim funkcji. Oznacza, ze funkcja nie ma zwracac
uzytecznej wartosci.

\begin{lstlisting}[style=moonchunk,caption={Funkcja bez wartosci zwrotnej}]
function logTitle(title: string): void {
  print(title);
  return;
}
\end{lstlisting}

\section{Deklaracja typu jawnie i niejawnie}
W wielu miejscach MoonChunk pozwala pominac jawny typ i oprzec sie na wartosci po prawej
stronie przypisania.

\begin{lstlisting}[style=moonchunk,caption={Deklaracja z typem niejawnym}]
let title = "MoonChunk";
\end{lstlisting}

Mozesz tez podac typ jawnie:

\begin{lstlisting}[style=moonchunk,caption={Deklaracja z typem jawnym}]
let title: string = "MoonChunk";
\end{lstlisting}

Jawne typy sa szczegolnie przydatne wtedy, gdy:

\begin{itemize}
  \item chcesz jasno opisac zamiar,
  \item zalezy Ci na czytelnosci,
  \item pracujesz z danymi z zewnatrz,
  \item deklarujesz \mc{let} bez wartosci poczatkowej.
\end{itemize}

\section{Zgodnosc typu przy przypisaniu}
Jesli deklarujesz typ, interpreter pilnuje zgodnosci. Przyklad poprawny:

\begin{lstlisting}[style=moonchunk,caption={Zgodne przypisanie typu}]
let count: int = 5;
count = 8;
\end{lstlisting}

Przyklad niepoprawny:

\begin{lstlisting}[style=moonchunk,caption={Niezgodne przypisanie typu}]
let count: int = 5;
count = "osiem";
\end{lstlisting}

Tutaj MoonChunk powinien zglosic blad typu.

\section{Globalne zmienne w \mc{env}}
Poza zwyklymi zmiennymi lokalnymi MoonChunk wspiera tez globalne deklaracje umieszczane
w bloku \mc{env}. To przydatne dla wartosci wspolnych dla wielu miejsc programu.

\begin{lstlisting}[style=moonchunk,caption={Globalne wartosci w bloku env}]
env {
  global siteName: string = "MoonChunk Portal";
  global pageSize: int = 10;
};
\end{lstlisting}

Wazna zasada: global deklarowany w \mc{env} powinien byc od razu zainicjalizowany.
Dzieki temu jest przewidywalny i gotowy do uzycia wszedzie tam, gdzie ma byc widoczny.

\section{Praktyczne wskazowki dotyczace typow}
\begin{itemize}
  \item Uzywaj \mc{int}, gdy wynik ma byc liczba calkowita.
  \item Uzywaj \mc{double} lub \mc{float}, gdy spodziewasz sie ulamkow.
  \item Uzywaj \mc{bool} do warunkow, a nie liczb zastępujacych prawde i falsz.
  \item Uzywaj \mc{object} lub \mc{dict} do grupowania nazwanych danych.
  \item Uzywaj \mc{array} do list i kolekcji.
  \item Po \mc{unknown} i \mc{any} siegaj dopiero wtedy, gdy zwykle typy nie wystarczaja.
\end{itemize}

\section{Najczestsze bledy poczatkujacych}
\begin{itemize}
  \item traktowanie \mc{const} jak \mc{let},
  \item odczytywanie \mc{let} przed inicjalizacja,
  \item mieszanie typu liczbowego z tekstem bez swiadomego rzutowania,
  \item wpisywanie wartosci logicznej tam, gdzie oczekiwana jest liczba,
  \item kopiowanie tej samej zmiennej w to samo miejsce zamiast wykonania zwyklego przypisania.
\end{itemize}

W nastepnym rozdziale przejdziemy od samych typow do wyrazen, czyli do tego,
jak MoonChunk oblicza wartosci i laczy je ze soba.
